\chapter{Génese Dependente I: A Noção de Eu é Ignorância}

\lettrine{A}{\space singularidade da abordagem budista} reside \emph{no anatta} -- a compreensão do não-eu. O estilo particular de reflexão em estruturas como as Quatro Nobres Verdades e o \emph{paticcasamuppada} transforma a forma de pensar, passando da visão egocêntrica -- do «dever» e do «eu» como absolutos -- para \emph{o anatta} -- o não-eu.

O problema reside no facto de que «não-eu» soa a aniquilação, não é verdade? E o que assusta as pessoas no budismo é que «não-eu» e «ausência de alma» soam como uma posição absoluta que se tem de assumir como budista. Pessoas que odeiam Deus e nutrem ressentimento pelo cristianismo podem tornar-se budistas porque guardam rancor contra Deus, a alma, o pecado e a culpa. Elas querem realmente que o budismo seja uma espécie de filosofia ateísta e uma rejeição total de toda a experiência cristã. Mas não é isso que ele é. O budismo não é ateísta nem niilista. O Buda foi muito cuidadoso em evitar tais posições extremas.

Em vez disso, o seu ensinamento é uma psicologia construída com muita habilidade e cuidado. O seu objectivo é ajudar-nos a ver através e a libertar-nos de todos os apegos habituais --- atitudes nascidas da ignorância, do medo e do desejo --- que criam esta sensação ilusória de um eu. Assim, há já 2530 anos que o budismo tem conseguido sobreviver e manter a sua pureza. E isso deve-se ao facto de a sua abordagem ser muito clara. Existe uma Sangha que vive sob a disciplina do Vinaya, e há o ensinamento do Dhamma.

Se praticarmos isto da maneira certa, podemos realmente começar a ver o sofrimento e a miséria que criamos por causa dessas ilusões sobre nós mesmos. Mas não estamos a tentar criar uma ilusão de que não existe nenhum eu. O objectivo não é passar da ilusão do eu para a ilusão de que «não existe eu», mas sim investigar, contemplar e ter a percepção até que a verdade inefável seja compreendida, cada um de nós por si mesmo.

Cada um de nós tem a sua experiência única -- não estamos a viver exactamente as mesmas coisas. Temos memórias, experiências, tendências e hábitos diferentes. E, no entanto, relacionamos sempre estas infinitas variedades com o Dhamma,
para que não estejamos apenas a fazer interpretações totalmente subjectivas. Aplicamos os ensinamentos do Dhamma à nossa experiência para podermos comunicá-la e compreendê-la num contexto mais amplo do que o da subjectividade pessoal.

Muitas vezes, as pessoas desviam-se na prática quando as suas experiências religiosas são interpretadas de forma demasiado subjectiva. Não são apresentadas de uma forma que possa ser comunicada. Tornam-se experiências pessoais únicas, em vez de realizações universais. Tal como aconteceu com o gnosticismo cristão, acaba-se por ter todo o tipo de interpretações subjectivas muito estranhas da experiência mística. Cada gnóstico tinha a sua própria maneira de falar, por isso a Igreja Católica Romana da época disse: «Isto é loucura!» e baniu tudo. Mas o Buda estabeleceu toda uma forma de pensar e de expressar os ensinamentos que é exactamente a mesma hoje. Não mudamos nem distorcemos tudo para se adequar à nossa experiência pessoal. Medimos a nossa experiência com o ensinamento porque os ensinamentos são tão habilmente elaborados que abrangem tudo.

Na contemplação do \emph{paticcasamuppada,} estamos a chegar a um consenso sobre como os seus termos se relacionam com a experiência contemplativa. Quando lemos \emph{o paticcasamuppada} pela primeira vez\emph{,} não percebemos nada. «A ignorância condiciona as formações kármicas; as formações kármicas condicionam a consciência, etc.» E então? O que é que isso significa? Imagina-se que deve ser algo muito profundo e que provavelmente é preciso uma vida inteira a estudar Pali para o compreender. Por isso, tende-se a ignorá-lo.

Nos círculos budistas, as Quatro Nobres Verdades podem ser ignoradas. «Ah, sim -- budismo básico. Sim\ldots{} agora vamos passar para o verdadeiro budismo Madhyamika avançado!» Ou: «O que disse Dogen?» Ou: «Milarepa é absolutamente fascinante, não é?» E pensas: «Sofrimento, Origem, Cessação e Caminho: sim, sabemos isso, agora vamos ao que realmente interessa.

Assim, as Quatro Nobres Verdades tendem a ser crenças superficiais. As pessoas não as investigam nem as utilizam porque os ensinamentos em si não são interessantes, pois não? Sofrimento, Origem, Cessação e Caminho; não é um ensinamento inspirador porque é um ensinamento para a prática -- não um ensinamento para inspirar. E é por isso que o utilizamos: porque essa forma específica de pensar e contemplar é, psicologicamente, de grande valor.

Com isto, podemos começar a compreender aquilo que nunca vimos ou compreendemos antes. Ao seguir esta forma de prática, estás, na verdade,
a desenvolver a sua mente e inteligência de uma forma que é muito autónoma. Mesmo nas formações mais avançadas, as pessoas não treinam realmente as suas mentes nesta forma específica de reflexão e contemplação. Pensar racionalmente é altamente valorizado. Mas para compreender o que é a racionalidade como uma função da mente, é preciso reflectir sobre a natureza da mente. O que está realmente a acontecer? Do que se trata tudo isto? E, claro, estas são as questões da existência, não são? As questões existenciais: «Por que nasci?» «A vida tem algum sentido?» «O que acontece quando morro?» «Do que se trata tudo isto?; ;Será que não tem sentido -- apenas um acidente cósmico?; ;Está relacionado com algo além de si mesmo ou é apenas algo que acontece e depois pronto -- é o fim?»

Temos grandes dificuldades em relacionar o sentido da vida com qualquer coisa real para além do mundo material. Assim, o materialismo torna-se a realidade para nós. Quando exploramos o espaço, é sempre no plano material. Queremos subir em foguetões e levar os nossos corpos até à Lua porque, de acordo com a visão materialista, é isso que é real. O materialismo ocidental carece de subtileza e requinte: leva-nos a um nível de consciência muito grosseiro, onde a realidade é o objecto material bruto e as emoções são descartadas como não sendo reais porque são subjectivas. Não se pode andar por aí a medir emoções com instrumentos electrónicos.

Mas as emoções, claro, são \emph{muito} reais para nós individualmente -- o que estamos a sentir é mais importante para nós do que o que está num relógio digital. Os nossos medos, desejos, amores, ódios e aspirações são o que realmente torna as nossas vidas felizes ou miseráveis. E, no entanto, estes podem ser descartados no materialismo moderno, num mundo baseado no prazer sensual, na riqueza material e no pensamento racional, de modo que a vida espiritual, para muitas pessoas, parece ser apenas uma ilusão. Não se pode medi-la com um computador nem examiná-la com instrumentos electrónicos.

No entanto, na civilização europeia pré-científica, o mundo espiritual era o mundo real. Como achas que construíram as catedrais? E toda esta arte surgiu de um verdadeiro sentido de aspiração espiritual, do ser humano ligado a algo para além do mundo material. A verdade espiritual é algo que cada um deve compreender individualmente. A verdade é a auto-realização, a subjectividade suprema. E o Buda leva a subjectividade até ao centro do universo, o ponto imóvel e silencioso, onde o sujeito não é um sujeito \emph{pessoal}. Esse ponto imóvel não é de ninguém nem de nada.

Na meditação, estás a caminhar nessa direcção. Estás a libertar-te de todos esses apegos às condições mutáveis do mundo material, do plano emocional, do plano intelectual, do plano simbólico e do plano astral. Tudo isso é abandonado para que possas alcançar o ponto de quietude, o silêncio. Esse desapego não é uma aniquilação ou uma rejeição, mas dá-te a perspectiva necessária para compreender o todo. Não podes compreender o todo estando na periferia, onde apenas és levado pelo turbilhão.

Ser levado pelo turbilhão na periferia significa que estás perdido no apego a todas as coisas que estão a girar à tua volta. Chama-se \emph{samsara.} Não tens capacidade para parar e observar porque estás simplesmente preso neste movimento circular.

O objectivo da meditação, nesta abordagem das Quatro Nobres Verdades e \emph{do paticcasamuppada,} é parar o rodopiar da mente. Na verdade, permaneces na quietude, não como um ataque contra o mundo condicionado, mas para o ver em perspectiva. Não o estás a aniquilar, nem a criticá-lo, nem a tentar escapar dele de forma alguma através da aversão ou do medo. Mas estás a chegar ao centro, ao ponto de quietude onde podes vê-lo tal como ele é e conhecê-lo, sem mais te assustares ou te iludires com ele. E fazemos isto dentro dos limites da nossa experiência pessoal. Assim, podemos dizer: «Cada um por si.» Porque é assim que parece quando estamos aqui sentados. E, no entanto, esse ponto de quietude não está na mente, não está no corpo; é aqui que se torna inefável. O ponto de quietude não é um ponto dentro do cérebro. No entanto, estás a perceber esse silêncio universal, essa quietude, essa unidade, onde todo o resto é um reflexo e é visto em perspectiva. E a personalidade, o kamma, as diferenças, as variedades e todas estas coisas já não nos iludem, porque já não nos agarramos a elas.

À medida que examinamos a mente cada vez mais, à medida que reflectimos e contemplamos sobre ela e aprendemos com ela, todos começamos a perceber a quietude da mente, que está sempre presente, mas que, na maioria das pessoas, nem sequer é notada. Isto porque a vida do \emph{samsara} é tão agitada, tão frenética, que nos deixa a girar. Embora o ponto de quietude esteja sempre aqui, nunca é visto até que se aproveite a ocasião para permanecer na quietude, em vez de andar à volta na circunferência.

Não que a quietude seja algo a que nos devamos apegar! Não estamos a tentar tornar-nos pessoas que estão quietas -- apenas sentadas aqui em quietude, sem sentir nada. Sei que alguns de vocês vêm aqui e criam um mundo pessoal que podem habitar durante a hora de meditação. Mas não é assim que
de sair do sofrimento; esse mundo subjectivo e pessoal depende muito de as coisas serem de uma determinada maneira. É tão frágil e tão efémero que é destruído com a mais leve perturbação. O mundo refinado da tranquilidade é tão encantador, tão pacífico -- e então alguém ressona! É repugnante, não é, ser tirado desses belos estados de tranquilidade por funções corporais grosseiras.

Mas a quietude não é tranquilidade. Não é necessário que estejamos tranquilos, mas há quietude quando conseguimos confiar em permanecer no silêncio, em vez de seguir as nossas tendências compulsivas. Todos tendemos a pensar que temos de estar a fazer alguma coisa; estamos tão condicionados a fazer coisas que até a meditação se torna uma actividade compulsiva em que estamos envolvidos. «Desenvolver isto\ldots{} desenvolve aquilo\ldots{} Tenho de desenvolver \emph{o} meu \emph{samadhi} e tenho de desenvolver os \emph{jhanas}.» Não se vem aqui apenas para se sentar; vem-se aqui para \emph{se desenvolver!} É assim que pensamos! Sentimo-nos culpados se não estivermos a fazer nada, a progredir, a desenvolver-nos, a chegar a algum lado. E, no entanto, ser capaz de entrar aqui e sentar-se em quietude não é uma coisa muito fácil de fazer, pois não? É muito mais fácil fazer grandes projectos de desenvolvimento de meditação, planos quinquenais e assim por diante. No entanto, acabamos sempre no ponto de quietude: a pensar nas coisas tal como elas são.

À medida que se compreende cada vez mais, pode haver um desapego do desejo de evoluir e tornar-se alguma coisa. E à medida que a mente se liberta de todo esse desejo de se tornar e obter algo, de alcançar algo, então a Verdade começa a revelar-se. Está sempre presente, aqui e agora. É uma questão de simplesmente ser capaz de estar aberto e sensível para que a Verdade se revele. Não é algo que se revela a partir do exterior. A Verdade está sempre presente, mas não a vemos se estivermos presos à ideia de conquistas, de «eu» ter de fazer algo, de «eu» ter de obter algo.

Por isso, o Buda lançou este ataque directo ao «eu e ao meu». É a única coisa que te está a bloquear. O obstáculo é o apego a uma visão do eu. Se simplesmente enxergares além dessa visão do eu, se a deixares ir, então compreenderás o resto. Não precisas de conhecer todas as outras fórmulas esotéricas elaboradas nem nada disso. Não precisas de te perder infinitamente na complexidade se simplesmente largares a visão ignorante de «eu sou».

Percebe isso, e conhece e compreende o caminho do desapego, do não apego. Então a Verdade revela-se onde quer que estejas, o tempo todo. Mas até o fazeres, estarás sempre preso a criar problemas e complicações.

\emph{«Avijjapaccaya sankhara; sankharapaccaya vinnanam; vinnanapacaya namarupam; namarupapaccaya salayatanam; salayatanapaccaya phassa; phassapaccaya vedana; vedanapaccaya tanha; tanhapaccaya upadanam; uapadanapaccaya bhavo; bhavapaccaya jati; jatipaccaya jaramaranam-soka-parideve-dukkha-domanassupayasa.»}\footnote{Esta é a formulação, na língua Pali das escrituras, do ensinamento sobre a Originação Dependente (paticcasamuppada). A tradução encontra-se na Introdução; o cerne da sua mensagem está contido no texto em português que se segue.}

Tudo isto significa que, se continuares a insistir em apegar-te às ilusões de um eu, à ganância, ao ódio e à ilusão, tudo o que alguma vez irás obter no final é velhice, doença e morte, luto, tristeza, desespero e angústia. É tudo o que irás obter para o resto da tua vida: uma perspectiva bastante enfadonha, não é!

Mas pode libertar-se disso aqui e agora, através do Compreensão Correta, vendo as coisas da maneira certa. Pode haver o conhecimento da Verdade, no qual já não se é iludido pelas aparências ou hábitos, ou pelas condições à nossa volta.

\photoPage{image-012-anagarika-washing.jpg}{Anagārika a lavar a louça}{Foto by Jeff Pick}


